\documentclass[11pt]{article}

\usepackage[a4paper,margin=2.4cm]{geometry}
\usepackage{booktabs}
\usepackage{array}
\usepackage{amsmath}
\usepackage{xcolor}
\usepackage[colorlinks=true,linkcolor=black,urlcolor=blue]{hyperref}
\usepackage{enumitem}
\setlist{nosep,leftmargin=1.4em}
\usepackage{titlesec}
\titlespacing*{\section}{0pt}{1.1ex}{0.6ex}
\titlespacing*{\subsection}{0pt}{0.9ex}{0.4ex}
\setlength{\parindent}{0pt}
\setlength{\parskip}{0.5ex}

\newcommand{\code}[1]{\texttt{\small #1}}

\title{\vspace{-1.4cm}\textbf{NextGenC --- Gridded SOC maps of Finland}}
\author{}
\date{}

\begin{document}
\maketitle
\vspace{-2.2em}

Procedural note for the kriged SOC map products that accompany the NextGenC
report. Companion to \code{build\_soc\_matrices.R} (point matrices) and
\code{build\_soc\_maps.R} (this product).

\section{What this produces}

For each of the 6 models (SP1, TP2, TP3, Yasso07, Yasso15, Yasso20), the
posterior-predictive SOC at the $\sim$447 calibration plots is interpolated to
a regular grid covering Finland, \textbf{one layer per year, 1985--2024
(40 years)}. Two multiband GeoTIFFs per model, plus a multi-model
\textbf{ensemble} pair (\textbf{14 files total}):

\begin{center}
\begin{tabular}{@{}lcl@{}}
\toprule
\textbf{File} & \textbf{Bands} & \textbf{Content} \\
\midrule
\code{<MODEL>\_SOC\_mean.tif} & 40 & Kriged mean total SOC (tC/ha), band $n$ = year $1984{+}n$ \\
\code{<MODEL>\_SOC\_sd.tif}   & 40 & Total uncertainty SD (tC/ha), same grid/year as mean band \\
\code{ENSEMBLE\_SOC\_mean.tif} & 40 & Six-model ensemble mean SOC (tC/ha) \\
\code{ENSEMBLE\_SOC\_sd.tif}   & 40 & Ensemble total SD (within + between-model + spatial) \\
\bottomrule
\end{tabular}
\end{center}

Band names carry the year (e.g.\ \code{y1985}), written into GeoTIFF metadata so
labels travel with the file. Band $n$ of every SD stack is the uncertainty for
band $n$ of the matching mean stack (same year, grid, and land mask).

\section{Grid and storage}

\begin{itemize}
\item \textbf{CRS:} ETRS89-LAEA Europe (EPSG:3035), metres; equal-area. The grid is
snapped to the EEA reference grid (cell edges on multiples of the 2\,km resolution
from the LAEA origin). Kriging is performed natively in this projection; source
inputs (plot coordinates, boundary, peat/water masks) are reprojected from
ETRS-TM35FIN (EPSG:3067).
\item \textbf{Resolution:} 2\,km. Grid $330\times589$ ($\sim$194\,k cells, of which $\sim$44\,\% land).
\item \textbf{Mask:} clipped to the Finland boundary from \code{maakunta\_sf.rds}.
\item \textbf{Format:} GeoTIFF, DEFLATE-compressed, float32.
\item \textbf{Footprint:} $\sim$0.12\,GB for all 560 maps (14 files).
\end{itemize}

The 2\,km choice is a deliberate compromise: the $\sim$447 plots have a
\textbf{mean spacing of $\sim$27.5\,km}, so 2\,km already oversamples the
information content --- finer grids render smoother but add no real spatial
detail. Disk budget (1--2\,GB) is not the binding constraint; data density is.

\section{Method}

Kriging is done \textbf{in log space} and back-transformed, because SOC is
positive and right-skewed and HIKET's error model is multiplicative --- this
keeps the map uncertainty multiplicative and consistent with the calibration.
Per model $\times$ year:

\begin{enumerate}
\item \textbf{Inputs.} Plot mean and SD come from the posterior-predictive
  bundles (the same \code{soc\_mean}\,/\,\code{soc\_sd} columns used by
  \code{build\_soc\_matrices.R}); plot coordinates (\code{x\_ETRS},
  \code{y\_ETRS}) come from \code{Data/model\_inputs/site\_raw.csv}, joined on
  \code{plot\_id}.
\item \textbf{Krige the mean.} Ordinary kriging of $\log(\texttt{soc\_mean})$
  over the grid $\rightarrow$ interpolated mean surface plus the \textbf{kriging
  prediction variance} $v_{\text{krige}}$ (small near plots, large in
  data-sparse areas, e.g.\ Lapland).
\item \textbf{Krige the model uncertainty.} Ordinary kriging of
  $\log(\texttt{soc\_sd})$ $\rightarrow$ a surface of the model's own
  posterior-predictive uncertainty.
\item \textbf{Combine (option C).} Total uncertainty fuses both sources,
  \[
    \text{sd}_{\text{total}} \;=\; \sqrt{\,v_{\text{krige}} \;+\; \text{kriged\_modelSD}^{2}\,}
  \]
  (combined in log space, then back-transformed). This propagates \textbf{model
  uncertainty} \emph{and} \textbf{spatial-coverage uncertainty} into a single SD
  layer --- neither source is silently dropped.
\item \textbf{Back-transform \& mask.} Exponentiate to tC/ha, clip to the land
  mask, write the two stacks.
\end{enumerate}

Variogram: \textbf{one pooled model per model}, reused for all 40 years (the
spatial structure of SOC is essentially static year-to-year; this is stable and
avoids per-year auto-fit failures). \code{set.seed(2025)} as elsewhere in HIKET.

\paragraph{Peat \& water exclusion.} The maps target \textbf{mineral-soil forest
SOC}. Peat/water calibration plots (\code{site\_raw} soil codes \code{Tu}/\code{Ve})
are dropped from the kriging \emph{input} so they do not pull the interpolated
surface, and grid cells are masked to NoData where they are \textbf{$>$50\,\%
peatland} or \textbf{$>$50\,\% water}. Peat: Luke MS-NFI site-main-class raster
(\code{paatyyppi}, 16\,m, EPSG:3067; spruce/pine mire + treeless peatland),
aggregated to the 2\,km grid as a peat fraction \emph{of the whole cell} (water in
the denominator, so lake-dominated cells do not read as peat). This site-type
definition captures drained peatland forest (\emph{turvekangas}), which land-cover
products miss ($\sim$10\,\% of land masked). Water: SYKE Ranta10 lake polygons
(the sea is outside the land polygon). See \code{Data/Peat\_extension/} and
\code{Data/Water\_extension/}. In the thumbnail figures peat renders \textbf{black}
and water/sea \textbf{white}.

\section{Ensemble (six-model) product}

The ensemble is formed \emph{on the grid}, after each model has been kriged ---
i.e.\ the per-model mean and total-SD stacks above are combined cell-by-cell,
equal-weighted ($w_i = 1/6$). With $M=6$ model mean surfaces $m_i$ and total-SD
surfaces $s_i$ (each $s_i$ already folds in that model's kriging variance and
posterior uncertainty, \S3 step~4):

\[
  m_{\text{ens}} = \frac{1}{M}\sum_i m_i,
  \qquad
  s_{\text{ens}} = \sqrt{\;
    \underbrace{\tfrac{1}{M}\textstyle\sum_i s_i^{2}}_{\text{within-model}}
    \;+\;
    \underbrace{\tfrac{1}{M-1}\textstyle\sum_i (m_i - m_{\text{ens}})^{2}}_{\text{between-model (structural)}}
  \;}.
\]

This is the law of total variance: the ensemble deviation fuses \textbf{(i)}
each model's own predictive uncertainty (which already carries \textbf{(iii)}
the spatial/kriging uncertainty) with \textbf{(ii)} the \emph{structural spread}
between models --- the among-model disagreement that is the central HIKET
question. The between-model term is the scientifically interesting one: where it
dominates, the map is uncertain because the models \emph{disagree}, not because
the data are sparse. Combination is done in the same (log/back-transformed)
space as the per-model SDs for consistency.

\textbf{Note:} the between-model term is also informative on its own; if a
structural-disagreement map is wanted later, export $\sqrt{\tfrac{1}{M-1}\sum_i
(m_i-m_{\text{ens}})^2}$ as a separate stack --- one extra line in the script.

\section{How to run}

\begin{verbatim}
# from the project root
Rscript Reporting/NextgenC_report/build_soc_maps.R
\end{verbatim}

Reads the same \code{bundles <- c(...)} RUN\_ID map as
\code{build\_soc\_matrices.R}. Writes the 12 GeoTIFFs into
\code{Reporting/NextgenC\_report/}.

\textbf{Dependencies:} \code{terra}, \code{gstat}, \code{sf} (and \code{sp}).
Posterior-predictive bundles must be present in
\code{Calibration\_real\_data\_transient/runs/}, and
\code{Data/model\_inputs/\{site\_raw.csv, maakunta\_sf.rds\}} must exist.

\section{Maintenance notes}

\begin{itemize}
\item \textbf{TP3 RUN\_ID is provisional.} As in \code{build\_soc\_matrices.R},
  the TP3 bundle is the pre-fix (Euler) predictive. Swap the TP3 RUN\_ID to the
  exact-integrator re-calibration when it is synced from Puhti, then re-run.
\item \textbf{Adding years / plots.} No code change --- the script reads year
  columns and plot rows from the bundles directly.
\item \textbf{Changing resolution.} Edit the single \code{RES\_M <- 2000}
  constant. Storage scales as $1/\text{res}^2$: 1\,km $\approx$ 0.4\,GB, 5\,km
  $\approx$ 0.02\,GB (all 560 maps, compressed).
\item \textbf{Reading the output.} \code{terra::rast("SP1\_SOC\_mean.tif")} gives
  a 40-layer SpatRaster; layer names are the years. Same for the \code{\_sd} file.
\end{itemize}

\end{document}
